% ============================================================
% Chapter 2: Literature Review
% CT-MUSIQ Undergraduate Thesis
% ============================================================

\chapter{Literature Review}
\label{chap:literature_review}

\section{Introduction to the Review}
\label{sec:lit_intro}

The problem of automated perceptual quality assessment for low-dose CT images sits at the
intersection of several active research domains: CT reconstruction and denoising, image
quality assessment theory and practice, deep learning with limited data, and Vision
Transformer architecture design. This chapter provides a structured review of the prior
art relevant to each of these domains, beginning with the clinical and technical context
of CT imaging and progressing through the evolution of IQA methods from hand-crafted
statistical approaches to modern neural network architectures. The chapter concludes with
a precise identification of the specific gap that CT-MUSIQ is designed to address.

The review is organised into seven thematic sections: (1) the clinical imperative and
evolution of CT reconstruction; (2) the dichotomy between subjective evaluation and
full-reference metrics; (3) ontological flaws in traditional no-reference IQA applied to
CT; (4) the deep learning revolution in NR-IQA and the feature extraction bottleneck;
(5) evaluation deficits and standardisation efforts; (6) Vision Transformers, multi-scale
processing, and the MUSIQ architecture; and (7) emerging frontiers in vision-language
model-based IQA. The final section synthesises the identified gaps and positions the
contributions of this thesis.


\section{The Clinical Imperative and Evolution of Computed Tomography Reconstruction}
\label{sec:lit_ct_reconstruction}

The expansion of versatile medical imaging modalities ranging from ultrasound and magnetic
resonance imaging (MRI) to positron emission tomography (PET) has fundamentally
revolutionised modern clinical diagnostics, yet it is computed tomography (CT) that
nowadays accounts for an enormous majority of large-scale, cross-sectional healthcare data
generation \cite{zhao2023multitask}. While conventional normal-dose CT (NDCT) provides
unparalleled, high-resolution anatomical visualisation, its diagnostic utility is
fundamentally connected with ionising X-ray radiation. This radiation raises significant
health concerns regarding radiation-driven genetic mutations and future cancer-related
risks. Guided strictly by the ALARA (As Low As Reasonably Achievable) principle, the
medical physics community has aggressively adopted Low-Dose CT (LDCT) protocols. The
critical efficacy of this adaptation has been heavily validated in expansive observational
studies that include trials for lung cancer screening and long-term patient tracking
\cite{cao2023observational}. Dose reduction is especially important in highly sensitive
demographics, such as pediatric populations undergoing ultra-low-dose abdominal scans or
cardiac CT evaluations, where optimising the signal-to-noise ratio (SNR) against radiation
risk is a delicate physiological necessity \cite{cho2024pediatric}.

Mathematically, image reconstruction is a typical inverse problem where the goal is to
recover an unknown signal $\mathbf{x}$ from a corrupted measurement
$\mathbf{y} = A(\mathbf{x}) + \epsilon$, with $A$ representing the forward corruption
process \cite{niu2023score}. In LDCT, reducing the incident X-ray photons inevitably
yields photon starvation at the detector array, which mathematically expands as a complex,
non-stationary ``Poisson + Gaussian'' degradation model; the quantum noise obeys a Poisson
distribution driven by photon count fluctuations, while the electronic noise of the data
acquisition system is characterised by a signal-independent Gaussian variance
\cite{zhang2024review}. Early analytical reconstruction methods, most notably Filtered
Back Projection (FBP), dominate clinical practice due to their computational efficiency
in calculating the analytical left inverse of the system matrix. However, FBP relies on
idealised linear assumptions that fail catastrophically under the high-noise,
photon-starved conditions of LDCT, resulting in major streak artefacts and extensive
noise amplification.

To overcome these computational limitations, Model-Based Iterative Reconstruction (MBIR)
and Hybrid Iterative Reconstruction (IR) algorithms were developed. These formulate the
reconstruction as a highly regularised objective function combining data fidelity with
complex structural priors, typically optimised using first-order methods or the
Alternating Direction Method of Multipliers (ADMM) \cite{goodfellow2016deep}. Although
IR techniques successfully incorporate statistical noise modelling (e.g., Markov Random
Field priors) to minimise mean squared errors, they are often computationally expensive
and frequently interfere with the inherent Noise Power Spectrum (NPS) and Modulation
Transfer Function (MTF) of the CT system \cite{xun2025bibliometric}. Additionally,
iterative reconstruction algorithms have been shown to drastically shift the visual
texture of the images, producing a non-natural, ``plastic'' or ``blotchy'' aesthetic that
significantly degrades the diagnostic conspicuity of low-contrast lesions
\cite{mileto2019iterative}.

Therefore, the industry has aggressively pivoted toward Artificial Intelligence and Deep
Learning Image Reconstruction (DLIR) algorithms, such as the Residual Encoder-Decoder CNN
(RED-CNN) or dual-domain frameworks, which map degraded low-dose sinograms or image
patches directly to new targets, significantly outperforming IR in preserving natural
noise textures and structural sharpness \cite{koetzier2023dlir}. Recently, DLIR has
evolved beyond basic feed-forward networks to embrace advanced generative approaches,
including the deployment of Score-Based Generative Models (SGMs) and sophisticated
diffusion predictive models like the Contextual Error-Modulated Generalised Diffusion
Model (CoreDiff) \cite{gao2024corediff}. Pushing the boundaries of physical modelling,
state-of-the-art frameworks such as the Noise-Inspired Diffusion Model (NEED) have
explicitly tuned the reverse diffusion Markov chain to mimic the specific noise
characteristics of CT, utilising a modified Poisson diffusion process to align the
generative pattern directly with the true degradation physics of pre-log LDCT projections
\cite{gao2025need}. Simultaneously, at the hardware level, the development of
Photon-Counting Detector CT (PCD-CT) is revolutionising multi-energy imaging by measuring
individual incident X-ray photons without electronic noise interference, demanding entirely
new generations of reconstruction networks \cite{willemink2013iterative}.

The proliferation of DLIR methods, each generating large quantities of reconstructed
images with varying quality characteristics, creates a pressing demand for automated,
reproducible quality metrics to replace manual radiologist evaluation for benchmarking
purposes. Furthermore, comparing FBP and advanced IR techniques explicitly reveals that
while iterative methods drastically improve standard signal-to-noise ratios, the
aggressive smoothing due to their algorithmic regularisation can reduce the diagnosis of
minute pathologies compared to the unprocessed raw data \cite{smits2007minor}. This
highlights why a perceptual, rather than purely signal-theoretic, quality metric is
essential for clinical relevance.


\section{The Dichotomy of Image Quality Assessment: Subjective Evaluation versus Full-Reference Metrics}
\label{sec:lit_iqa_dichotomy}

In the highly regulated medical domain, image quality is not defined by graphical fidelity
alone but rather by the measurable contribution to diagnostic accuracy. The ultimate gold
standard for Image Quality Assessment (IQA) remains the subjective, diagnosis-based
evaluation performed by highly trained, board-certified radiologists who assess images for
critical clinical criteria, including lesion visibility, subtle structural preservation,
and the suppression of hallucinated artefacts \cite{nirupama2025}. To measure the
reliability of these human annotations, statistical indicators such as the Intraclass
Correlation Coefficient (ICC) are extensively implemented to quantify both inter-observer
agreement (consistency across different radiologists) and intra-observer reliability
(repeatability of a single radiologist's assessments over time) \cite{herath2025}.

However, subjective IQA is limited by unavoidable practical bottlenecks: it is
labour-intensive, financially exorbitant, and inherently non-scalable for evaluating
large groups of computational outputs. Early computational attempts to automate this
assessment relied on complex mathematical model observers, such as the Channelised
Hotelling Observer (CHO) and the Channelised Support Vector Machine (CSVM). These models
utilise Singular Value Decomposition (SVD) and Principal Component Analysis (PCA) to
perform dimensionality reduction on specific Signal-Known-Exactly/Background-Known-Exactly
(SKE/BKE) detection tasks \cite{pouget2023channelised}.

To automate the assessment of broader, generalised image degradations without
task-specific requirements, objective Full-Reference IQA (FR-IQA) metrics were eagerly
adopted from the computer vision domain. These algorithms compute the degradation of a
test image through direct mathematical comparisons against a perfectly aligned reference
counterpart, utilising metrics such as Peak Signal-to-Noise Ratio (PSNR), Structural
Similarity Index Measure (SSIM), Feature Similarity (FSIM), Visual Signal-to-Noise Ratio
(VSNR), Information Fidelity Criterion (IFC), and Visual Information Fidelity (VIF)
\cite{yan2018deeplesion}. To properly demonstrate the true clinical utility of these
metrics and the underlying deep learning denoising networks, researchers construct
sophisticated virtual imaging simulations based directly on patient raw data; by inserting
simulated lesions and noise into the projection domain, they can accurately measure the
oversampled Edge Spread Function (ESF) and in-plane MTF at various clinical contrast
levels (e.g., $-10$, $-20$, $-50$~HU) \cite{zhou2024virtual}.

Despite demonstrating high mathematical correlation in strictly controlled phantom studies,
FR-IQA possesses an insurmountable ontological flaw for genuine clinical application.
Acquiring a perfectly aligned, artefact-free, high-dose NDCT reference image for every
corresponding LDCT scan is physically, practically, and ethically impossible due to strict
radiation safety constraints and the inevitability of patient motion between sequential
scans \cite{athar2019comprehensive}. This fundamental limitation provides the primary
rationale for developing no-reference (NR-IQA) methods that can operate without any
paired reference image.

The transition from FR-IQA to NR-IQA, while operationally necessary, introduces
significant representational challenges. Without a reference, the model must infer quality
degradation from the statistics of the corrupted image alone. For natural photographs,
this has been achieved with considerable success by exploiting the statistical regularity
of high-quality natural scenes. However, medical images do not share these statistical
properties, as explored in the following section.


\section{Ontological Flaws of Traditional No-Reference IQA for CT}
\label{sec:lit_nr_iqa_failures}

To eliminate the strict reliance on reference data, No-Reference IQA (NR-IQA), frequently
referred to as ``blind IQA'', was systematically developed to evaluate image quality based
solely on statistical features extracted from the degraded image itself, transforming the
assessment problem into a classification or regression task \cite{wang2006modern}. The
most established traditional NR-IQA algorithms, such as the Blind/Referenceless Image
Spatial Quality Evaluator (BRISQUE) and the Natural Image Quality Evaluator (NIQE),
operate on the foundation of Natural Scene Statistics (NSS) \cite{hu2024niqe}. These
algorithms extract handcrafted statistical features by calculating local Mean Subtracted
Contrast Normalised (MSCN) coefficients; they rely on the fundamental assumption that
high-quality, distortion-free photographic images generally exhibit a predictable,
heavy-tailed Gaussian distribution, whereas distorted images deviate from this norm,
allowing models to fit deviations using a Generalised Gaussian Distribution (GGD) or an
Asymmetric GGD (AGGD) \cite{sheikh2005nss}.

While this framework is mathematically elegant and effective for assessing generic
photographic degradations like JPEG compression or optical blur, the direct application of
NSS-based NR-IQA to medical CT results in catastrophic performance degradation
\cite{athar2019comprehensive}. Natural photographic images are shaped by broader spatial
frequencies generated by reflected visible light, whereas clinical CT images represent
highly specific, biophysical X-ray absorption coefficients precisely mapped to Hounsfield
Units (HU). The complex, non-stationary background noise power spectrum and the
spatially variant, macro-structural streak artefacts characteristic of photon-starved
LDCT fundamentally violate the core assumptions of natural scene broadband models,
rendering these classical NR-IQA metrics completely unreliable for predicting true
diagnostic usefulness.

The LDCTIQAC 2023 challenge results corroborate this quantitatively: BRISQUE achieves an
Aggregate score of only 2.1219 (PLCC~$= 0.75$, SROCC~$= 0.7863$, KROCC~$= 0.5856$)
on the leaderboard \cite{lee2025ldctiqac}, demonstrating the severe underperformance of
classical statistics-based approaches on calibrated medical images. The SNR-based baseline
similarly achieves only 2.4026, confirming that even physically motivated signal metrics
do not capture the perceptual quality criteria applied by radiologists.

Furthermore, this catastrophic breakdown of traditional statistical blind IQA is not
isolated to CT. Analogous failures have been extensively documented in Magnetic Resonance
(MR) imaging, where classical NR metrics are entirely incapable of accurately quantifying
modality-specific degradations such as ghosting, aliasing, or spatial displacement of
water and fat molecules \cite{stepien2023mri}. These cross-modality failures underscore
the need for dedicated, domain-adapted NR-IQA systems trained with clinical expert labels,
rather than unsupervised statistical models derived from natural image databases.


\section{Deep Learning-Driven NR-IQA and the Feature Extraction Bottleneck}
\label{sec:lit_dl_iqa}

The representational limitations of handcrafted statistical features led to the massive
integration of deep learning (DL) into the NR-IQA paradigm, fundamentally transforming
the field from manual statistical modelling to data-driven representation learning
\cite{ahmed2024iqanrtl}. Early deep learning IQA solutions aggressively took advantage
of pre-trained CNN backbones, such as VGG, ResNet, DenseNet, and Inception architectures,
by eliminating their classification heads and utilising transfer learning to map structured,
high-dimensional feature representations to continuous subjective quality scores via
fully connected regression layers.

To further enhance feature extraction and prevent the loss of critical low-level visual
information due to deep pooling layers, sophisticated architectures like the staircase
structure were developed to systematically integrate features from intermediate network
layers, ensuring the model leverages both low-level textural noise signatures and
high-level semantic anatomical context simultaneously \cite{hosu2020koniq}. To establish
explicit medical domain priors into the network's learning trajectory, researchers
developed highly specialised hybrid architectures. A paramount example involves fusing
Sparse Representation encoding with psychophysical Just Noticeable Distortion (JND)
models. By explicitly calculating patch-level JND maps based on luminance adaptation and
texture masking gradients, these networks selectively pool features only from visually
critical, distortion-sensitive anatomical regions, thereby mimicking the localised
attention of the Human Visual System (HVS) \cite{shen2025jnd}.

Advanced deep learning frameworks have also been engineered for specific, highly complex
imaging models, such as Cone-Beam CT (CBCT), where sparse-view projections are repeatedly
optimised using U-Net-based restoration operators driven by time-embedding modules to
suppress severe undersampling artefacts \cite{greenspan2023miccai_cbct}. For standard
LDCT, CNN-based IQA has been applied using encoder-decoder networks with transfer learning
from 2D to 3D analysis, demonstrating that cross-dimensional knowledge transfer can
partially address the limited data problem \cite{shan2018transfer}.

Despite these significant architectural improvements, pure CNN-based NR-IQA algorithms
suffer from a critical geometric limitation. Standard CNNs fundamentally rely on
localised convolution kernels and stacked spatial pooling to incrementally expand their
receptive fields. As a direct consequence, they invariably reduce and destroy fine-grained,
high-frequency quantum noise information in deeper layers, while simultaneously failing
to establish the global, long-range spatial dependencies required to accurately evaluate
macro-structural streak artefacts that pass through the entire anatomical cross-section.
Furthermore, training these deep CNNs on diverse clinical data from multi-institutional
sources introduces severe domain shift challenges, driving the necessity for advanced
active learning frameworks and incremental domain adaptation strategies to effectively
scale model generalisability without requiring exhaustive, dense manual delineations
\cite{greenspan2023miccai_domain}.


\section{Evaluation Deficits and Standardisation}
\label{sec:lit_standardisation}

The historical progression of deep learning-based LDCT denoising and automated IQA has
been severely slowed by systemic, widespread flaws in experimental evaluation protocols
across the scientific literature. A comprehensive literature analysis reveals significant
deficits: the implementation of insufficient objective metrics that have little to no
correlation with actual clinical diagnostic value, the highly unfair selection of
hyperparameters optimised specifically to negatively impact competing reference models,
and an overwhelming reliance on small, homogeneous, privately held datasets that prevent
robust model generalisation \cite{mccollough2020ldct}. To comprehensively address this
critical shortcoming in algorithmic standardisation and benchmarking, researchers utilise
highly sophisticated physical simulations to generate low-dose estimation data from
routine-dose scans; these simulations precisely model the Poisson statistics of X-ray
photons and the non-Gaussian components of detector electronic noise, creating accurate
datasets at precise dose reductions (e.g., 40\%, 20\%, and 10\% simulated doses)
\cite{zabic2013simulation}.

Establishing standardised, high-fidelity datasets is not simply an academic exercise for
IQA evaluation, but an absolute prerequisite for ensuring the reliability of
next-generation automated medical analysis tasks, such as the precise, deep-learning-driven
segmentation of COVID-19 infections and pulmonary lesions from highly degraded CT volumes
\cite{sori2021dfdnet}.

The most significant advancement in standardising perceptual assessment was the
\textbf{Low-Dose Computed Tomography Perceptual Image Quality Assessment Grand Challenge
(LDCTIQAC)}, held in association with MICCAI 2023 and subsequently documented in Lee
et al.\ \cite{lee2025ldctiqac}. The challenge curators released an unprecedented,
massive open-source dataset comprising 1,000 abdominal CT images sourced from
multi-institutional scanners, meticulously annotated by expert radiologists using a
strictly calibrated 5-point Likert grading scale to embed clinical perception directly
into the ground truth labels.

The challenge attracted multiple competitive teams, whose leaderboard results provide the
primary benchmark for comparison in this thesis:

\begin{itemize}
    \item \textbf{agaldran} (Aggregate~$= 2.7427$): The top-performing submission,
    combining Swin Transformer and ResNet-50 backbones, with multi-head model calibration
    using the approach described in \cite{galdran2022calibration}.

    \item \textbf{RPI\_AXIS} (Aggregate~$= 2.6843$): Second-place submission.

    \item \textbf{CHILL@UK} (Aggregate~$= 2.6719$): Third-place submission.

    \item \textbf{FeatureNet} (Aggregate~$= 2.6550$): Fourth-place submission.

    \item \textbf{SNR baseline} (Aggregate~$= 2.4026$): Physics-motivated baseline
    using signal-to-noise ratio measurements.

    \item \textbf{BRISQUE} (Aggregate~$= 2.1219$): Classical NSS-based blind IQA
    applied directly without adaptation.
\end{itemize}

To further reduce the severe risk of catastrophic overfitting when training massive deep
neural networks on relatively small, specialised medical datasets, advanced regularisation
and data augmentation strategies, such as the unpredictable, convex interpolation of
image-label pairs known as MixUp, are being extensively investigated and integrated into
state-of-the-art training pipelines to force the networks to learn robust, linear
transitions between distinct quality degradation classes \cite{galdran2022calibration}.


\section{Vision Transformers, Multi-Scale Processing, and the MUSIQ Architecture}
\label{sec:lit_vit_musiq}

To absolutely outperform the localised mathematical inductive biases of standard CNNs, the
computer vision community has recently executed a massive architectural pivot toward the
Vision Transformer (ViT) approach, which visualises an image as a flattened sequence of
tokenised pixel patches and employs multi-head self-attention mechanisms to dynamically
compute complex correlations between all tokens across the entire image \cite{sheikh2005nss}.
This convolution-free mechanism is uniquely suited for LDCT analysis, as it establishes a
global receptive field from the very first layer, enabling the network to simultaneously
correlate distant, low-frequency streak artefacts and preserve overall anatomical structure
without downsampling.

\subsection{Vision Transformer Foundations}
\label{subsec:vit_foundations}

The foundational Vision Transformer (ViT) model, introduced by Dosovitskiy et al.\ in
2020, demonstrated that a pure self-attention architecture without convolutional inductive
biases could achieve competitive or superior performance on large-scale image recognition
benchmarks when pretrained on sufficiently large datasets. The core ViT operation
tokenises an image of dimensions $H \times W$ into $N = HW/P^2$ non-overlapping patches
of size $P \times P$, projects each patch to a $d$-dimensional embedding, prepends a
learnable \texttt{[CLS]} token, adds position encodings, and processes the resulting
sequence through $L$ Transformer encoder blocks consisting of Multi-Head Self-Attention
(MHSA) and position-wise feed-forward networks (FFN).

The MHSA operation for a query matrix $Q$, key matrix $K$, and value matrix $V$ is:

\begin{equation}
    \text{Attention}(Q, K, V) = \text{softmax}\!\left(\frac{QK^\top}{\sqrt{d_k}}\right) V
\end{equation}

where $d_k = d / H_{\text{attn}}$ is the per-head dimension. The global attention operation
allows every patch to attend to every other patch simultaneously, enabling the model to
capture long-range structural dependencies that CNN architectures require many layers to
approximate.

However, applying standard ViT architectures to NR-IQA introduces a significant technical
barrier known as the \textbf{input resolution paradox}: rigid ViTs require fixed-resolution
inputs to maintain constant sequence lengths for their learnable position embeddings,
meaning native high-resolution medical slices must be globally resized. This approach
fundamentally destroys the high-frequency quantum noise signatures critical for quality
assessment, while rigid crops eliminate the global anatomical context necessary to evaluate
macro-artefacts.

\subsection{The Swin Transformer and Hierarchical Attention}
\label{subsec:swin}

The Swin Transformer introduced a hierarchical design using shifted-window attention that
partially addresses the resolution challenge by computing attention within local windows
and using shifted window partitioning to introduce cross-window connections. The
Swin-T backbone achieves competitive performance on dense prediction tasks while
maintaining linear computational complexity with respect to image size. In the context of
LDCT denoising, the Swin Transformer-based Encoder-Decoder Network (STEDNet) \cite{zhu2023stednet}
dramatically improves noise suppression by utilising shifted-window mechanisms and skip
connections for global feature fusion, constrained by a hybrid L1 and MS-SSIM loss
function. Similarly, the Hybrid CNN-Transformer codec network (HCformer) \cite{luo2025hcformer}
utilises Window/Sliding-Window Multi-head Self-Attention (W/SW-MSA) to efficiently
capture information interaction across the entire image. Furthermore, sophisticated
convolution-free networks like TED-net \cite{fan2021interpretability}, multi-domain
optimisation networks (MDST) that integrate Swin Transformers directly into both
projection and image domain sub-networks for sparse-view CT \cite{li2023mdst}, and
Cross-Network Contextual (CNA) modules \cite{zhang2021convergent} are setting new
benchmarks.

\subsection{The MUSIQ Architecture}
\label{subsec:musiq}

The Multi-Scale Image Quality Transformer (MUSIQ) \cite{ke2021musiq}, proposed by Ke
et al.\ at ICCV 2021, represents a powerful algorithmic solution to the resolution
paradox specifically in the context of NR-IQA. Rather than resizing images to a single
fixed resolution, MUSIQ extracts patches from a multi-scale spatial pyramid consisting of
two or more target resolutions, concatenates all patches from all scales into a single
sequence, and processes this variable-length sequence through a standard Transformer
encoder.

To handle the variable sequence length arising from different scales, MUSIQ introduces a
\textbf{hash-based 2D spatial embedding mechanism} that encodes each patch's scale index,
row index, and column index as three independent learned embedding lookups, summed to form
the positional encoding:

\begin{equation}
    \text{pos}_{i} = \mathbf{E}_{\text{scale}}(s_i) + \mathbf{E}_{\text{row}}(r_i) + \mathbf{E}_{\text{col}}(c_i)
\end{equation}

This factored design allows MUSIQ to process images of arbitrary resolution without
interpolation-induced information loss, explicitly encoding the scale identity and spatial
coordinates of each patch. MUSIQ achieved state-of-the-art performance on multiple
natural IQA benchmarks (KonIQ-10K \cite{hosu2020koniq}, SPAQ, etc.) at the time of its
publication, demonstrating the superiority of the multi-scale Transformer approach over
CNN-based methods for quality-aware feature extraction.

\subsection{Limitations of MUSIQ for Medical CT}
\label{subsec:musiq_limitations}

Despite its architectural elegance and strong performance on natural image benchmarks,
the vanilla MUSIQ architecture is not directly applicable to the medical CT IQA domain
for several reasons:

\textbf{Natural image preprocessing assumptions.} MUSIQ normalises input images using
ImageNet mean and standard deviation statistics (designed for RGB natural photographs),
which is inappropriate for CT images calibrated in Hounsfield Units. The resulting
mismatch between training distribution and CT pixel value distributions would cause the
pretrained patch embeddings to operate outside their trained domain.

\textbf{Absence of clinical windowing.} Radiologists who annotated the LDCTIQAC dataset
applied brain soft-tissue windowing before scoring. A model that does not apply identical
windowing before feature extraction is processing a fundamentally different image than the
one being scored, creating a systematic misalignment between the annotation and the model's
perceptual input.

\textbf{No scale-consistency regularisation.} The original MUSIQ does not enforce any
agreement between per-scale quality predictions, potentially allowing the model to rely
on scale-specific artefacts rather than learning truly scale-invariant quality features.
On a small medical dataset of 700 training samples, this risk of scale-specific
overfitting is significantly elevated.

\textbf{Small dataset regime.} MUSIQ was designed for and evaluated on large natural image
IQA datasets (KonIQ-10K has 10,000 images; SPAQ has 11,125 images). Adapting it to 700
training samples requires additional regularisation strategies beyond those present in
the original architecture.

These limitations motivate the specific architectural adaptations introduced in CT-MUSIQ,
as detailed in Chapter~\ref{chap:methodology}.


\section{Emerging Frontiers: Vision-Language Models in IQA}
\label{sec:lit_vlm}

Currently, the highest frontier of medical image quality assessment involves combining
Large Language Models (LLMs) and Multi-modal Large Language Models (MLLMs) to replicate
the cognitive, step-by-step diagnostic reasoning of clinical radiologists \cite{wu2024mllm}.
State-of-the-art architectures, such as IQAGPT, leverage expansive foundational pathways
(e.g., PaLM, LLaMA) fine-tuned with highly specific, autoregressive language modelling
targets directly on CT-IQA datasets \cite{chen2024iqagpt}. By passing visual tokens
extracted from image encoders through complex cross-modal attention mechanisms,
Vision-Language Models (VLMs), including highly advanced frameworks like LLaVA
\cite{liu2023llava} and InternLM-XComposer-VL, are capable of generating remarkably
accurate visual descriptions and reasoning paths.

Through the implementation of strict mental and physical prompting strategies, these models
seamlessly convert abstract numerical quality ratings into highly semantic, comprehensive
textual descriptions (e.g., explicitly detailing the presence of ``photon starvation
streaks blocking the hepatic blood vessel'' or ``optimal preservation of low-contrast
lesions in the abdominal soft-tissue window''), generating complete, explainable diagnostic
quality assessment reports along with absolute numerical scores.

Architectures such as InstructBLIP \cite{dai2023instructblip} and CLIP-based models
\cite{radford2021clip} have demonstrated strong zero-shot and few-shot transfer capabilities
that are beginning to be explored for medical imaging tasks. The incorporation of
domain-specific medical knowledge into these large pre-trained models through instruction
tuning represents a promising direction for reducing the dependency on large labelled
medical datasets.

While these foundational VLMs offer unprecedented zero-shot explanations and represent a
monumental leap toward transparent medical AI, they are currently limited by massive
computational overhead, critical inference latency, and the persistent risk of
hallucinating clinical metadata. Consequently, their immediate, practical deployment as
low-latency, high-throughput objective quality scoring substitutes within standard,
real-time hospital PACS pipelines remains highly limited, requiring further optimisation
before widespread clinical adoption. Furthermore, fully understanding the underlying
optimisation principles of deep transformers, specifically phenomena such as ``grokking''
where networks exhibit sudden, delayed leaps in generalisation performance long after
achieving perfect training accuracy, remains a critical frontier in theoretical AI
research \cite{elhage2021mathematical}.


\section{Synthesis and Identification of the Research Gap}
\label{sec:lit_gap}

The comprehensive framework of the deep learning literature clearly suggests that automated
NR-IQA algorithms must strictly simulate the precise perceptual characteristics of expert
radiologists to be considered clinically acceptable. The following key conclusions emerge
from the review:

\textbf{Classical NR-IQA is insufficient.} NSS-based methods like BRISQUE and NIQE fail
catastrophically on CT images due to the mismatch between natural image statistics and
the Hounsfield Unit calibration domain. LDCTIQAC challenge results provide quantitative
confirmation of this failure with BRISQUE achieving only Aggregate~$= 2.1219$.

\textbf{CNN-based methods are architecturally limited for CT IQA.} The localised
inductive biases of convolutional networks make them poorly suited for capturing the
global, long-range streak artefacts that dominate LDCT quality degradations.

\textbf{Vanilla ViT architectures face the resolution paradox.} Fixed sequence length
requirements force interpolation of high-resolution CT slices, destroying the
high-frequency noise signatures critical for perceptual quality assessment.

\textbf{MUSIQ resolves the resolution paradox but was not designed for CT.} The
multi-scale hash-based tokenisation of MUSIQ is architecturally ideal for LDCT IQA,
but the original model assumes natural RGB image statistics and lacks the clinical
preprocessing and regularisation adaptations necessary for the CT domain.

\textbf{The specific research gap is:} a dedicated MUSIQ-derived architecture that
enforces clinical CT windowing as a preprocessing stage, incorporates scale-consistency
regularisation to prevent overfitting on small medical datasets, and is systematically
evaluated against the LDCTIQAC 2023 grand challenge benchmark.

CT-MUSIQ is proposed to directly address this gap. By synthesising cutting-edge
Transformer architectures with the physical constraints of clinical neuroradiology, the
CT-MUSIQ framework delivers a computationally efficient, clinically grounded, and
empirically validated automated LDCT quality assessment engine. The specific architectural
and training innovations that constitute CT-MUSIQ are described in full in the following
chapter.


\section{Transfer Learning and Domain Adaptation in Medical Imaging}
\label{sec:lit_transfer}

A critical enabling technology for CT-MUSIQ is the transfer of representations learned
on natural image datasets to the medical imaging domain. This section reviews the
principles and practice of transfer learning in medical imaging contexts, with emphasis
on the specific challenges and strategies relevant to quality assessment on small
annotated datasets.

\subsection{Fundamentals of Transfer Learning for Medical Imaging}

Transfer learning in deep neural networks leverages the observation that features learned
by networks trained on large, diverse datasets (such as ImageNet with 1.28 million
labelled images across 1,000 categories) are partially general-purpose: lower-layer
features such as edge detectors, Gabor-like filters, and colour gradients transfer broadly
across visual domains, while higher-layer features become increasingly task-specific and
less transferable \cite{goodfellow2016deep}. In the context of medical imaging, this
observation motivates a two-phase training strategy: initialising a network with
weights pretrained on natural images, then fine-tuning on the (smaller) medical dataset.

The degree to which natural image pretraining transfers to CT imaging depends on the
architectural depth at which domain mismatch occurs. For shallow CNNs, the first few
convolutional layers learn general low-level features (edges, textures) that are relevant
to CT images; deeper layers trained for natural image classification may not transfer
as effectively due to the fundamentally different statistical structure of CT images.
Vision Transformers, however, exhibit a different transfer pattern: the self-attention
mechanism learned on natural images captures relational patterns between patches that
generalise more broadly than the rigid spatial inductive biases of convolution,
potentially enabling better medical image representations from the same pretrained
starting point.

Several studies have systematically investigated this transfer property for medical
imaging. Pretraining on ImageNet consistently outperforms random initialisation as a
starting point for medical image classification, detection, and segmentation tasks, even
when the target domain (CT, MRI, histopathology) appears visually dissimilar to natural
photographs. The benefit is especially pronounced for small medical datasets where training
from random initialisation would overfit severely. In the CT IQA context, this
justifies the use of ViT-B/32 pretrained on ImageNet as the encoder backbone for
CT-MUSIQ, despite the obvious domain gap between natural RGB photographs and
single-channel greyscale CT images calibrated in Hounsfield Units.

\subsection{Grayscale-to-RGB Adaptation Strategies}

A practical challenge in transferring ImageNet-pretrained models to single-channel CT
images is the three-channel (RGB) assumption embedded in the pretrained weights. Three
principal strategies exist for handling this mismatch:

\textbf{Strategy 1: Channel repetition.} The greyscale image is replicated to all three
channels, producing an RGB image where R = G = B = pixel value. This is the simplest
approach and preserves the pretrained weights intact. The trade-off is that the model
processes redundant information in all three channels, potentially wasting capacity on
the trivial inter-channel correlations. However, for the patch embedding layer, channel
repetition ensures that the linear combination of the three-channel input through the
pretrained convolutional filter weights produces a meaningful activation, since the filter
was designed to respond to spatially correlated patterns across channels.

\textbf{Strategy 2: Channel average initialisation.} The pretrained RGB patch embedding
weights are collapsed to a single channel by averaging across the input channel dimension,
and the single-channel architecture is used. This reduces the parameter count of the
patch embedding layer by a factor of three but requires re-initialisation of the patch
embedding, losing the pretrained weights.

\textbf{Strategy 3: Domain-specific pretraining.} The entire model is pretrained from
scratch on large unlabelled medical image datasets using self-supervised objectives
(e.g., masked image modelling as in MAE, or contrastive learning as in MoCo) before
fine-tuning on the annotated CT quality labels. This approach is computationally expensive
but potentially produces better initialised features for the CT domain.

CT-MUSIQ adopts Strategy~1 (channel repetition) as the most practical approach given the
small 700-image training set and limited compute budget. The simplicity of this approach
is supported by the empirical observation that even with channel repetition, the model
achieves competitive quality assessment performance, suggesting that the pretrained
ViT-B/32 features provide strong initialisation regardless of the greyscale-to-RGB
conversion approach.

\subsection{Fine-Tuning Stability on Small Medical Datasets}

Fine-tuning large pretrained models on small datasets presents significant stability
challenges. The key risks are: (1) catastrophic forgetting of pretrained representations,
where aggressive fine-tuning overwrites the useful pretrained features; (2) overfitting
to the limited training distribution; and (3) gradient instability due to the mismatch
between the pretrained feature scale and the randomly initialised task head gradients.

Several strategies have emerged to address these challenges in the medical imaging
literature. Layer-wise learning rate decay, where earlier layers are fine-tuned with
smaller learning rates than later layers, has been shown to preserve lower-level
pretrained features while allowing higher-level task-specific adaptation. Prompt-based
tuning approaches, which freeze the entire pretrained backbone and only train a small
set of learnable prompt tokens injected into the input sequence, provide an extreme
form of parameter-efficient adaptation that has shown promise for few-shot medical image
analysis.

CT-MUSIQ addresses fine-tuning stability through the two-stage training protocol described
in Chapter~\ref{chap:methodology}: Stage~1 warms up the randomly initialised
domain-specific components (prediction heads, hash positional encoding) with the encoder
frozen, preventing their noisy initial gradients from corrupting the pretrained attention
patterns. Stage~2 then unfolds gradually with linear warm-up followed by cosine
annealing, allowing the encoder to adapt smoothly to the CT domain while the heads
continue to refine their quality mapping.


\section{Regularisation Strategies for Small Dataset Deep Learning}
\label{sec:lit_regularisation}

The challenge of training deep learning models on small datasets like the 700-image
CT-MUSIQ training set is a recurring theme in medical image analysis. This section
reviews the principal regularisation strategies relevant to this setting, providing
the motivation for the specific choices made in CT-MUSIQ.

\subsection{Data Augmentation for CT Images}

Data augmentation is the most widely used regularisation strategy for small medical
datasets, as it artificially increases the effective training set size by applying
label-preserving transformations to the original images. However, CT images impose
significantly stricter constraints on acceptable augmentations than natural photographs,
as discussed in Section~\ref{sec:augmentation}.

Beyond the mild augmentations used in CT-MUSIQ (flips, small rotations, brightness
jitter, Gaussian noise), more aggressive augmentation strategies such as CutMix, MixUp,
and random erasing have been explored for medical image analysis. MixUp creates
convex combinations of image-label pairs, which has been shown to improve calibration
and reduce overconfident predictions. However, for CT IQA, the quality scores of mixed
images do not follow a simple convex combination of the component scores (the perceptual
quality of a noise-blended CT image is not predictable from the component scores),
making MixUp semantically invalid without careful adaptation. The LDCTIQAC 2023 challenge
top submission (agaldran) employed MixUp with specific adaptations, as reported in
Galdran et al.\ \cite{galdran2022calibration}, demonstrating that well-calibrated
augmentation can improve performance on this dataset.

\subsection{Dropout and Weight Decay}

Dropout, originally introduced by Srivastava et al., is a widely used regularisation
technique that randomly zeroes a fraction of activations during training, preventing
co-adaptation of neurons and forcing more distributed feature representations. For
Vision Transformers fine-tuned on small datasets, reducing the dropout rate below the
standard 0.1 has been found to improve performance: the pretrained features already
provide a form of implicit regularisation, and excessive dropout can prevent the model
from fully leveraging the informative pretrained attention patterns.

Weight decay (L2 regularisation) applied through the decoupled AdamW optimiser provides
an additional regularisation force that pushes parameter magnitudes toward zero,
preventing overly large weights that generalise poorly. In CT-MUSIQ, weight decay is
set to 0.01, a value that provides meaningful regularisation without preventing
convergence.

\subsection{Scale-Consistency KL Regularisation as a Novel Form of Self-Supervision}

The scale-consistency KL divergence loss introduced in CT-MUSIQ can be understood as a
form of self-supervised consistency regularisation, conceptually related to the
consistency training paradigm in semi-supervised learning. In consistency training,
the model is trained to produce similar predictions for two different augmented versions
of the same input; this forces the model to learn representations that are invariant
to the specific augmentation applied.

In CT-MUSIQ, the ``two views'' of the same CT image are the coarse-scale patches
($224 \times 224$) and the fine-scale patches ($384 \times 384$). These two views
present the same underlying image content at different spatial resolutions, and the
scale-consistency KL loss enforces that the model's quality assessment is consistent
across these two views. This is semantically valid because the true diagnostic quality
of a CT image does not depend on the resolution at which it is evaluated---a streak
artefact visible at fine resolution is also present in the coarse-scale representation,
merely at lower spatial frequency.

From a statistical learning theory perspective, the KL regularisation can be interpreted
as an inductive bias toward scale-invariant quality representations. By penalising
per-scale predictions that differ substantially from the global prediction, the loss
encourages the model to learn features that are simultaneously visible and discriminative
at both resolutions. This is particularly valuable for CT IQA because different classes
of quality degradations manifest at different scales: global contrast reduction is
visible at both scales equally, streak artefacts are more conspicuous at fine scale,
while overall structural noise is better captured at coarse scale.


\section{Evaluation Metrics in IQA Research}
\label{sec:lit_metrics}

A comprehensive understanding of the evaluation metrics used in IQA research is
essential for interpreting the experimental results in this thesis. This section reviews
the theoretical properties of the three correlation metrics used in the LDCTIQAC 2023
challenge evaluation, with emphasis on what each metric measures and why their joint
assessment provides a more complete picture of quality prediction performance.

\subsection{Pearson Linear Correlation Coefficient}

The Pearson Linear Correlation Coefficient (PLCC) measures the strength of the linear
relationship between predicted and true quality scores. A PLCC of 1.0 indicates perfect
linear agreement; PLCC of 0 indicates no linear relationship; PLCC of -1 indicates
perfect negative linear correlation. The PLCC is sensitive to the scale and offset of
the predictions: a model that systematically overestimates scores by a constant will
have a high PLCC (if the linear relationship is preserved) but poor absolute accuracy.

In practice, PLCC rewards models that produce scores whose magnitude tracks the true
quality scale accurately. It is the most natural metric for evaluating the quality of
a regression model in the absolute sense, and is sensitive to the specific calibration
of the prediction scale. A low PLCC may indicate either that the model cannot distinguish
quality levels (poor rank ordering) or that its scores are poorly calibrated on the
absolute scale (even if the relative ordering is correct).

\subsection{Spearman Rank-Order Correlation Coefficient}

The Spearman Rank-Order Correlation Coefficient (SROCC) measures the agreement of the
rank orderings of predicted and true scores, independently of any linear or monotone
rescaling of the predictions. Formally, SROCC is the Pearson correlation of the rank
variables $R(\hat{q})$ and $R(q)$, where $R(\cdot)$ denotes the rank function.

SROCC is invariant to any monotone transformation of the predicted scores, making it
a robust metric that does not penalise systematic scale errors. This property makes
SROCC particularly appropriate as an evaluation metric for IQA, where the goal is to
correctly order images by quality rather than to produce scores on a specific absolute
scale. High SROCC indicates that the model can reliably order images from worst to best
quality, which is the core requirement for practical applications such as scan rejection
(reject scans below a threshold) and comparative benchmarking.

\subsection{Kendall Rank-Order Correlation Coefficient}

The Kendall Rank-Order Correlation Coefficient (KROCC) measures the proportion of
concordant pairs (pairs where the model's ordering agrees with the true ordering) minus
the proportion of discordant pairs, normalised by the total number of pairs. It is a
more granular measure of local ordering accuracy than SROCC: SROCC is determined by
the global rank transformation, while KROCC counts each pair independently.

KROCC is generally harder to maximise than SROCC and is more sensitive to local ordering
errors. A model with high SROCC but lower KROCC is one that gets the overall rank
ordering roughly right but makes more errors at the pair level---for instance, it
correctly identifies that group A images are higher quality than group B images, but
makes more errors in distinguishing within groups. For clinical applications where the
precise relative ordering of specific pairs of images matters (e.g., selecting the
best scan from a repeated acquisition), KROCC is the most directly relevant metric.

\subsection{Aggregate Score}

The LDCTIQAC 2023 challenge uses the sum of absolute values of all three metrics
(PLCC + SROCC + KROCC) as the primary ranking criterion, with a theoretical maximum
of 3.0. This composite metric rewards models that perform well across all three
aspects of quality prediction: absolute calibration (PLCC), global rank ordering
(SROCC), and local pair ordering (KROCC). A model that excels in only one or two of
these aspects will have a lower aggregate than one that performs well across all three,
making the aggregate a comprehensive measure of overall IQA performance.

The aggregate metric also provides a natural weighting: since all three metrics have
similar scales (all bounded in $[-1, 1]$), they contribute approximately equally to the
aggregate, preventing any single metric from dominating the evaluation. This is
particularly important when comparing methods with different architectures or training
objectives, some of which may be more naturally suited to maximising one metric over
the others.
